\section{Cached Field Arithmetic}
\label{sec:cache}

Arithmetic in a finite field is limited to a small set of results determined by the field's cardinality \(q\).
For addition and multiplication, this requires up to \(q^2\) entries in each cache.
The reciprocal cache is smaller, with at most \(q\) entries.
Division does not require a separate cache because it is implemented as multiplication by the reciprocal.

The implementation stores each cache in an
\texttt{std::unordered\_map<int64\_t, int64\_t>}.
Using integer keys avoids a custom hash function for the finite-field polynomial type.
Before an arithmetic result can be cached, each polynomial is converted to an integer index.
Listing~\ref{lst:index} shows the \texttt{index()} implementation.

\begin{lstlisting}[style=cpp,
  caption={Conversion from a finite-field polynomial to an integer index},
  label={lst:index}]
std::uint64_t index() const
{
    std::uint64_t result = 0;
    std::uint64_t mult = 1;

    for (std::size_t i = 0; i < K; i++)
    {
        result += value_.coefficient(i).representative() * mult;
        mult *= P;
    }

    return result;
}
\end{lstlisting}

The cache mapping key for an operation involving elements with indices \(i\) and \(j\) is
\[
k(i,j)=i+qj.
\]
The existing \texttt{from\_index()} method converts the resulting stored index back into a polynomial.
Reciprocal elements are stored as a one-to-one mapping between an element index and the index of its reciprocal.
The cache entries are loaded once when the field is initialized.

The updated implementation is located in the
\texttt{algebra-finite-fields-group2} repository:
\url{https://github.com/JacobNance142/algebra-finite-fields-group2}.
The custom \texttt{FiniteField} class was included in the
\texttt{algebra-assignment} CMake configuration, allowing the determinant
and matrix code to use the cached operations without a public-interface
change.
